\section{Conclusion}

Multilingual chemistry retrieval needs evaluation data that makes the right
failure modes visible. In this paper, we presented a QAC generation and
validation pipeline for building such data, released it in retrieval-ready form,
and used it to evaluate dense multilingual retrievers beyond aggregate recall.
The central lesson is simple: a model that looks effective on average may still
prefer same-language shortcuts or chemically plausible but wrong documents over
the intended cross-language evidence.

The benchmark shows how to make those failures measurable. Language metadata
separates same-language from cross-language relevance; QAC provenance makes
generated queries auditable; human validation checks the LLM-based verifier; and
alias-graph hard negatives expose chemistry-specific confusions. These pieces
are useful beyond the particular patent-derived corpus used here: they provide a
template for building specialized multilingual retrieval evaluations in
high-trust technical domains.

For practitioners, the recommendation is to report robustness next to recall.
For benchmark builders, the recommendation is to start from the workflow failure
and preserve enough metadata to diagnose it. In multilingual chemistry search,
that means asking not only whether a relevant document was retrieved, but
whether the system found the intended evidence across language and chemical
confusability boundaries.
